% ! TeX root = ./main.tex
\section{Challenges of Telemetry–Driven Placement}
\label{sec:amr:chal}


\figchal

In principle, fine-grained telemetry should enable placement decisions to be grounded in observed
workload behavior.
However, our early attempts revealed two fundamental difficulties.
First, the telemetry itself was unreliable: metrics such as communication time showed poor
correlation with predictors like message volume (\cref{fig:amr:chal:regr}), making it impossible to
model baseline runtime.
Second, the magnitude of this variability also changed unpredictably as placement properties
were modified.
% Second, the artifacts driving this variability behaved unpredictably as placement properties were
% modified.
For example, optimizing for communication locality did not reliably improve end-to-end performance initially.
These observations forced us to treat placement not as a theoretical optimization problem, but as a
systems engineering task grounded in empirically observed runtime behavior.
The following challenges structure the empirical and algorithmic contributions of this work.


\parahead{Challenge 1: Cross-stack Performance Anomalies}
When initial telemetry proved inconsistent, showing poor correlation between workload metrics and
runtime (\cref{fig:amr:chal:regr}), our first challenge was establishing trust in the measurements.
In production environments, organizational and operational boundaries both constrain and scope
analyses.
Our full-stack access eliminated these boundaries: every deviation, across hardware, drivers, MPI,
or application layers, became our direct responsibility to diagnose and explain.
While this significantly increased diagnostic complexity, it also provided the crucial advantage of
unrestricted access to low-level tools like \texttt{perf}~\cite{Linux:LinuxPerf} and
eBPF~\cite{:eBPF}, which were necessary to establish the measurement fidelity underlying this
work.

\parahead{Challenge 2: Analyzing Fine-Grained Telemetry}
Diagnosing anomalies required analyzing large volumes of per-timestep, per-rank, and per-block
telemetry.
Anomalies originated on a small subset of ranks but propagated globally through communication,
making root cause identification nontrivial (see \cref{fig:amr:chal:vol}).
Tracing tools~\cite{Shend2006:TAU,Mey2012:Score-P} produced unstructured, high-volume output in
formats like OTF2~\cite{Mey2012:Score-P} unsuited for query-driven diagnosis.
Our analytics workflow evolved organically, beginning with standard tooling such as
TAU~\cite{Shend2006:TAU} and progressing towards custom query-driven workflows capable of surfacing
low-level system effects at scale.

\parahead{Challenge 3: Effective, Low-Overhead Placement}
Reliable telemetry exposed genuine load imbalance that could be targeted by placement, but
incorporating this information imposed strict new constraints.
Placement decisions had to be computed within tight time budgets, under 50\,ms per redistribution
for our target codes, to avoid measurable overhead on simulation runtime.
Even simplified placement objectives, such as minimizing makespan, are
NP-hard~\cite{Garey1979:Intractability}, and additional objectives further complicate optimization.
Placement mechanisms had to encode relevant tradeoffs while remaining fast and robust enough for
dynamic AMR workloads.

We address these challenges sequentially: \S\ref{sec:amr:tele} describes the tuning and analytics
infrastructure used to establish reliable telemetry, and \S\ref{sec:amr:design} presents placement
policies enabled by this exercise.